7\documentclass[parskip=full]{article}

%%%%%%%%%%%%%%%%%%%%%%%%%%%%%%%%%%%%%%%%%%%%%%%%%
% Basic includes
%%%%%%%%%%%%%%%%%%%%%%%%%%%%%%%%%%%%%%%%%%%%%%%%%
\usepackage{array}
\usepackage{amsmath}
\usepackage{amsfonts}
\usepackage{verbatim}
\usepackage{mathrsfs}
\usepackage{amssymb}
\usepackage{amsthm}
\usepackage{bbm}
\usepackage{bm}
%\usepackage[capposition=top]{floatrow}
\usepackage[dvipsnames]{xcolor}
\usepackage{cancel}
\usepackage{ulem}
\usepackage[labelfont=bf]{caption}
\usepackage{subcaption}
\usepackage{graphicx}
\usepackage{float}
\usepackage{natbib}
\usepackage{booktabs}
\usepackage{longtable}
\usepackage{multirow}
\usepackage{rotating}
\usepackage{pdflscape}
\usepackage[title]{appendix}
\usepackage{footmisc}
\usepackage{bigdelim}
\usepackage{tabularx}
\usepackage{xspace}
\newcommand{\A}{\ensuremath{\mathcal{A}}\xspace}
\newcommand{\B}{\ensuremath{\mathcal{B}}\xspace}
\newcommand\pa[1]{\ensuremath{\left(#1\right)}}
\usepackage{epigraph}
\usepackage{setspace}
%\onehalfspacing
\setstretch{1.5}

%%%%%%%%%%%%%%%%%%%%%%%%%%%%%%%%%%%%%%%%%%%%%%%%%%%
% Microtype
%%%%%%%%%%%%%%%%%%%%%%%%%%%%%%%%%%%%%%%%%%%%%%%%%%%
\usepackage[activate={true,nocompatibility},final,tracking=true,kerning=true,spacing=true,factor=1100,stretch=10,shrink=10]{microtype}
% activate={true,nocompatibility} - activate protrusion and expansion
% final - enable microtype; use ``draft'' to disable
%tracking=true, kerning=true, spacing=true - activate these techniques
% factor=1100 - add 10% to the protrusion amount (default is 1000)
% stretch=10, shrink=10 - reduce stretchability/shrinkability (default is 20/20)

%%%%%%%%%%%%%%%%%%%%%%%%%%%%%%%%%%%%%%%%%%%%%%%%%%%
% Hyperref setup
%%%%%%%%%%%%%%%%%%%%%%%%%%%%%%%%%%%%%%%%%%%%%%%%%%%
\usepackage{url}
\makeatletter
\g@addto@macro{\UrlBreaks}{\UrlOrds}
\makeatother
\usepackage{hyperref}
\hypersetup{
  colorlinks,%
  citecolor=black,%
  filecolor=black,%
  linkcolor=blue,%
  urlcolor=black
}

%%%%%%%%%%%%%%%%%%%%%%%%%%%%%%%%%%%%%%%%%%%%%%%%%%%
% Font setup
%%%%%%%%%%%%%%%%%%%%%%%%%%%%%%%%%%%%%%%%%%%%%%%%%%%
\usepackage{mathpazo} % math & rm
\linespread{1.5}        % Palatino needs more leading (space between lines)
\usepackage{parskip}
\usepackage[scaled]{helvet} % ss
\usepackage{courier} % tt
\usepackage{natbib}

\usepackage{filecontents}

\normalfont
\usepackage[T1]{fontenc}

\usepackage{titlesec}
\titleformat{\section}
{\normalfont\fontsize{12}{17}\bfseries}
{\thesection}
{1em}
{}

\titleformat{\subsection}
{\normalfont\fontsize{11}{16}\bfseries}
{\thesubsection}
{1em}
{}

%%%%%%%%%%%%%%%%%%%%%%%%%%%%%%%%%%%%%%%%%%%%%%%%%
% macros
%%%%%%%%%%%%%%%%%%%%%%%%%%%%%%%%%%%%%%%%%%%%%%%%%
\newcolumntype{C}[1]{>{\centering\let\newline\\\arraybackslash\hspace{0pt}}m{#1}}
\newcommand{\indicator}[1]{\mathbbm{1}_{\left\{ {#1} \right\} }}
\newcommand{\resultsloc}{../code/python/output/}
\newcommand{\subfigwidth}{0.475\textwidth}
\newcommand{\subfigwidthtwo}{0.375\textwidth}
\newcommand{\figwidth}{0.6\textwidth}
\DeclareMathOperator*{\argmax}{arg\,max}

\setlength\epigraphwidth{0.9\textwidth}
\setlength\epigraphrule{0pt}
\renewcommand{\epigraphflush}{center}
\usepackage{etoolbox}
\makeatletter
\patchcmd{\epigraph}{\@epitext{#1}}{\itshape\@epitext{#1}}{}{}
\makeatother

%%%%%%%%%%%%%%%%%%%%%%%%%%%%%%%%%%%%%%%%%%%%%%%%%
% Page and paragraph settings
%%%%%%%%%%%%%%%%%%%%%%%%%%%%%%%%%%%%%%%%%%%%%%%%%
\setlength{\parskip}{0.0pc} \setlength{\parindent}{0.25in}
\setlength{\topmargin}{-3pc}
\setlength{\textheight}{9in}
\setlength{\oddsidemargin}{0pc} \setlength{\evensidemargin}{0pc}
\setlength{\textwidth}{6.5in}
%\renewcommand{\footnotelayout}{\setstretch{1.2}}
\makeatletter
\@fpsep\textheight
\makeatother

%%%%%%%%%%%%%%%%%%%%%%%%%%%%%%%%%%%%%%%%%%%%%%%%%
% Title page
%%%%%%%%%%%%%%%%%%%%%%%%%%%%%%%%%%%%%%%%%%%%%%%%%
%%%%%%%%%%%%%%%%%%%%%%%%%%%%%%%%%%%%%%%%%%%%%%%%%
% Title page
%%%%%%%%%%%%%%%%%%%%%%%%%%%%%%%%%%%%%%%%%%%%%%%%%

\author{
  Shahar Rotberg\footnote{University of Toronto. Contact: \url{shahar.rotberg@mail.utoronto.ca}.} \and Lin Zhang\footnote{CMHC. Contact: \url{lzhang@cmhc-schl.gc.ca}. We would like to thank Joseph B. Steinberg for kindly allowing us to use his high-powered computer to compute the model and our simulations in this paper.}
}

\title{\vspace{-1.8cm}Rent Control: A Good Idea?}

\pagestyle{plain}

\begin{document}

\maketitle

\setstretch{1.275}
\begin{abstract}
\onehalfspacing
\noindent Does rent control improve housing affordability and societal welfare? To answer these questions, we use a general-equilibrium model with endogenous housing and rental markets, and with tenure duration choice. We discipline the model with data from Toronto, Canada where rent control is essentially universal, and then examine the consequences of rent deregulation. 
Deregulation would reduce renters' tenancy duration and increase homeownership. It would also reduce housing affordability, but benefit most households and raise welfare because deregulation would reduce rents early in life and increase them later in life, thus better smoothing consumption over households' lifecycles.\\[1.5ex]
% words=95

\noindent \textbf{JEL Classification:} E17; R31; R52. 

\noindent \textbf{Keywords:} General equilibrium, rent control, housing markets.
\end{abstract}
\setstretch{1.5}

\section{Introduction}

\section{Background: Rent control in Toronto, Canada}

\section{Related literature}

\section{Model}

The model economy features overlapping generations of finitely-lived households, a rental company, a housing construction company, and a government. Households in each cohort are heterogeneous in labor market productivity, housing tenure, years they have lived in the same rental, housing wealth, and financial wealth. Before retirement, households work in the labor market and after retirement they receive social security benefits from the government that depend on their average lifetime labor income. As in many other models in the housing literature, the interest rate, $r$, is exogenous \citep{sommeretal-2013,arslanetal-2015,sommersullivan-2017,kaplanetal-2019}; The economy is modelled as a small open economy. Homeowners pay proportional property taxes and maintenance costs each period. Housing and rental housing are supplied by competitive companies with convex costs. The government finances public goods, which do not enter households' utility, social security benefits to retired households, and lump sum transfers, by levying taxes on income and property.

\subsection{Demographics}\label{subsection:demographics}

Households are born at age $j=1$, retire from the labor force at age $j=R$, and live to $J$ years. Households do not face mortality risk, and after age, $J$, they die with certainty. Households do not care about their children, and thus in this set up, optimally households leave no bequests for their children. Newborn households begin their lives as assetless renters. Households discount the future at a rate of $\beta$.

Similarly to \citet{sommersullivan-2017}, households have the following preferences:
\begin{equation*}
    u(c,h,o,\iota)=\frac{\left\{c^{1-\gamma} \left[h(1+o\eta)\right]^{\gamma}\right\}^{1-\sigma}}{1-\sigma},
\end{equation*}

\noindent where $c$ represents non-housing consumption, $h$ denotes housing services consumption, $o=1$ represents a homeowner and $o=0$ represents a renter, $\eta$ represents the utility advantage to homeownership, $\gamma$ is the household's utility weight on the consumption of shelter services, and $\sigma$ is the household's constant relative risk aversion. Non-housing consumption is taxed at a rate $\tau_{c}$. 

\subsection{Endowments}

We denote financial wealth by $a\in \mathcal{A}=\mathbb{R}_+$, tenure by $o \in \mathcal{O}=\{0,1\}$, years lived in the same rental by $\iota \in \mathcal{I}=\mathbb{N}$, house size by $h\in \mathcal{H}=\mathbb{R}_{++}$, and mortgage debt by $m\in \mathcal{M}=\mathbb{R}_+$. Households are endowed with a stochastic labor market productivity, $e \in\mathcal{E} = \mathbb{R}_{++}$. $\Psi_{j}(a,o,\iota,h,m,e)$ is the measure of age-$j$ households over endowments, and we use $\mathcal{S}=\mathcal{A}\times\mathcal{O}\times\mathcal{I}\times\mathcal{H}\times\mathcal{M}\times\mathcal{E}$ to denote the household's state space. To simplify on notation, we mostly use $s$ to denote a household's state $(a,o,\iota,h,m,e)$.

Newborn households' labor market productivities are drawn from a conditional distribution \\ $\hat{F}(e^{\mbox{child}}|e^{\mbox{parents}})$, and then follow a Markov process $F(e'|e)$ until mandatory retirement. Labor market productivity also has a deterministic lifecycle component, $\zeta_j$, that increases with age; a household of age $j$ with labor market productivity $e$ supplies $e\zeta_j$ efficiency units of labor. Lastly, each household receives a lump sum transfer $T$ from the government each period, which we set to zero in the benchmark.

\subsection{Housing}\label{subsection:housing_choice}

Households can choose to rent or own a house, $o \in \mathcal{O}=\{0,1\}$, and the size of their unit, $h$. As is common in the literature, the smallest unit is $\underline{h}>0$ \citep{gervais,sommeretal-2013,sommersullivan-2017,kaplanetal-2019}. The purchase price of a house of size $h$ is $p h$. Because of rent control, the longer a renter stays in their unit, the cheaper their real rent becomes, and thus the cost of renting as function of the number of years a renter has lived in the same unit is $\frac{p^{r}}{(1+\kappa)^{\iota}} h$. $p_{r}$ is cost to rent, $\kappa$ determines the rate at which rent control reduces real rents over time, and $\iota$ is the number of years a renter has lived in the same unit. Buying and selling a house entail one-time proportional transaction costs, $\tau_{b} p h$ and $\tau_{s} p h$, respectively. Each period, housing is taxed at a rate $\tau_{p}$ and depreciates in value at a rate $\delta_h$. 

Before retirement, owners have access to mortgage debt, $m$, which is limited to $m \leq \lambda_{m}(ph) p h$, where $\lambda_{m}(ph)<1$, and depends on the value of the house. We set the interest rate on mortgage debt to $r_{m}=r+\Delta r$, where $r$ is the return on financial wealth, and $\Delta r$ is the interest premium paid on mortgage debt \citep{sommeretal-2013,sommersullivan-2017,kaplanetal-2019}. Following \citet{gervais}, mortgage debt is fully repaid at the beginning of each period. Because in Toronto households cannot deduct their mortgage interest or property tax payments from their taxable income, they cannot do so in the model either.

\subsection{Labor and capital incomes}\label{subsection:income}

Working-age households supply labor inelastically, and earn an exogenous wage $W$ for each unit of labor they supply. A working-age household whose age is $j$ and who has labor market productivity, $e$, earns labor income that is given by $y_{\ell,j}(e)=We\zeta_j(e)$. Instead of labor income, retired households receive social security transfers that depend on their expected labor income conditional on their labor productivity at retirement, $y_{R}(e) = B(e)$. The parameter $B(e)$ represents the fraction of the average labor income that a retiree with labor market productivity, $e$, receives. Households that retire with high productivities are more likely to have enjoyed high productivities throughout their lives, and thus receive larger social security transfers. Households also earn capital income $y_{k}(a)=ar$ on their financial wealth, $a$, where $r$ is the rate of return. We define taxable income as

\begin{equation*}\label{eq:taxable_labor_income}
  y^{x}_{j}(e,a) =y_{\ell,j}(e)+\nu y_{k}(a),
\end{equation*}

\noindent where, $\nu$ is the fraction of capital income that is taxable. Then, following \citet{hsv-2020}, we define income taxes to be progressive as follows:

\begin{equation}\label{eq:income_tax}
  y^{\tau}_{j}(e,a) = y^{x}_{j}(e,a)-\tau_{\ell}y^{x}_{j}(e,a)^{\psi},
\end{equation}

\noindent where $\tau_{\ell}$ determines the average level of taxation in the economy and $\psi$ measures the degree of progressivity in the tax system.

\subsection{Dynamic program}\label{subsection:household_optimization}
Given its state variables, a household chooses how much to save ($a'$), whether to own or rent ($o' \in \left\{0,1\right\}$), whether to remain in their unit one period longer ($\iota'$), and how much housing to consume ($h$), and how much mortgage debt to take on ($m'$). The dynamic program that represents this problem in a stationary equilibrium is

\begin{align}
  \label{eq:vf}
  V_{j}(s)=\max_{c,a',o',\iota',h',m'}\ \biggl\{u\left(c,h,o'\right)+
  \beta \mathbb{E}\biggl[V_{j+1}(s'|e)\biggl]\biggr\}
\end{align}
subject to
\begin{flalign}
  \begin{split} \label{eq:bc}
  & c(1+\tau_{c})+a'+(1-o')\frac{p^{r}}{(1+\kappa)^{\iota'-1}}h'+o'(ph'-m')+o\left(ph\left(\tau_{p}+\delta_{h}\right)+r_{m}m\right) + \\ & \quad \indicator{(o'=1\land o=0)\vee\left(o'=1 \land o=1 \land h' \neq h\right)}\tau_{b}ph' +\indicator{(o'=0\land o=1)\vee\left(o'=1 \land o=1 \land h'\neq h\right)}\tau_{s}ph= \\ & \quad  y_{\ell,j}(e)+y_{k}(a)-y^{\tau}_{j}(e,a)+a+o(ph-m)+ T   
  \end{split} \\
  & c \geq 0 \label{eq:consumption}\\
  & a' \geq 0 \label{eq:financial_wealth}\\
  & o' \in \{0,1\} \label{eq:tenure}\\
  & h' \geq \underline{h}  \label{eq:housing_renter}\\
  & \lambda_{m}(ph') p h'\geq m'\geq 0 \hspace{0.2cm} \mbox{if} \hspace{0.2cm} o'=1 \hspace{0.2cm} \mbox{and} \hspace{0.2cm} j\leq R\label{eq:mortgage_debt_owner}\\
  & m'=0 \hspace{0.2cm} \mbox{if} \hspace{0.2cm} o'=0 \hspace{0.2cm} \mbox{or} \hspace{0.2cm} j>R \label{eq:mortgage_debt_renter}
\end{flalign}

The household's utility is made of two components. The first is the flow utility from non-housing consumption and housing services consumption. The second is the household's expected continuation value based on its choices discounted by the household's discount factor, $\beta$.\footnote{We include labor market productivity shocks for retired households as well as non-retired households in our formulation of the dynamic program for notational brevity. Writing the dynamic program this way allows me to use a single equation to represent both retired and non-retired households' problems.} 

Equation \ref{eq:bc} represents the household's budget constraint. The left-hand side contains expenditures: consumption plus consumption tax; saving; the cost of renting a unit as a function of the years they have lived in the same unit; the purchase price of a house using a mortgage for a homeowner; property taxes, maintenance costs, and mortgage interest costs for a household starting the period as a homeowner; buying costs for a household starting the period as a renter and that chooses to become a homeowner or for a household starting the period as a homeowner that decides to purchase a different property; and selling costs for a household starting as an owner and that decides to sell and move to a different unit as an owner or as a renter. The right-hand side contains resources available to spend: after-tax labor and capital incomes; non-housing financial wealth; housing equity for a household that begins the period as an owner; and a lump sum transfer. 

Constraints \ref{eq:consumption} and \ref{eq:financial_wealth} restrict households to have positive consumption and savings. Constraint \ref{eq:tenure} says that the tenure choice can take the value $1$ for owners, or $0$ for renters. Constraints \ref{eq:housing_renter} requires households to consume at least $\underline{h}$ units of housing. Constraint \ref{eq:mortgage_debt_owner} requires working-age owners to have positive mortgage debt that cannot go above a fraction $\lambda_{m}(h')$ of their chosen house value. Constraint \ref{eq:mortgage_debt_renter} says that renters and retirees must have zero mortgage debt. I denote the policy functions for consumption, savings, tenure, years living in the same unit, housing, and mortgage debt by $c_{j}(s)$, $a'_{j}(s)$, $o'_{j}(s)$, $\iota'_{j}(s)$, $h'_{j}(s)$, and $m'_{j}(s)$ respectively.

\subsection{Rental companies}

We use the structure in \citet{garrigaeetal-2009}. For each $x=1$ units of rental housing, a rental company obtains rent, $p^{r}_{i}$, pays property taxes, $\tau_{p}$, and incurs depreciation costs, $\delta_{h}$. Where $i \in \{f,c\}$ represents the free rental market and controlled rental market, respectively. Note that in the controlled rental market $p^{r}_{i}$ is the average rent, $\overline{p}^{r}$, paid by a household living in a rental controlled unit. To supply a total of $X'_{i}$ units of rental housing, a rental company also incurs management costs equal to $c_{m,i}X'^{\theta}_{i}$, where $c_{m,i}$ determines the average cost to manage the stock of rental housing $X'_{i}$ and $\theta$ determines how quickly the marginal cost of supplying rental housing rises with the amount of rental housing supplied. Because the rent-to-price ratio in Toronto is quite low, it does not cover the flow of operating costs, which would imply a net loss to the rental company. To fix that, we add another parameter, $g_{p}$, which exogenously accounts for the real growth in house prices in Toronto, and thus guarantees that rental companies have incentives to supply rental housing. A rental company can, without friction, buy housing for the housing price, $p$, and each period a rental company decides how much rental housing, $X'_{i} \geq 0$, to rent out. Thus, a rental company has the following problem:

\begin{equation*}
V(X_{i})=\underset{X'_{i} \geq 0} {\text{$\mbox{max} $}}\left\{p^{r}_{i}X'_{i}-c_{m,i}X'^{\theta}_{i}-p X'_{i}+p(1-\delta_{h}-\tau_{p}+g_{p})X_{i}+\frac{1}{1+r}V(X'_{i})\right\}.
\end{equation*}

\noindent Combining the first order condition with respect to $X'_{i}$ with the Envelope condition gives the following equation for a rental price in market $i$ in the stationary equilibrium:

\begin{align}
&p^{r}_{i}(X'_{i})=c_{m,i}\theta X'^{\theta-1}_{i}+p-\left\{\frac{1-\delta_{h}-\tau_p+g_{p}}{1+r}\right\}p, \quad i \in \{f,c\}.\label{eq:rental_price}
\end{align}

\noindent Equation (\ref{eq:rental_price}) demonstrates that when $\theta>1$, the rental price in market $i$ must rise in order for the quantity of rental housing supplied to increase. When $\theta=1$, the rental market in market $i$ collapses to the traditional rental housing model in which the supply of rental housing is infinitely elastic. In our calibration we pick the value of $\theta$ so that the model generates a tax incidence between the rental companies and its renters that is aligned with estimates in the literature.

\subsection{Housing construction company}

Following \citet{rotberg-2022}, housing is supplied endogenously by a construction sector, which allows us to account for the impact of policy changes on the supply of housing, and thus housing prices and rents. The law of motion for housing is $H'=(1-\delta_h)H + I$, where $I$ denotes the construction sector's investment in new housing stock. For its production of $I$ new units of housing, the construction company receives payments equal to $pI$, but incurs an increasing marginal cost, as captured by the following production technology:

\begin{equation*}
&TC(I)=c_{1}\left(\frac{1}{\frac{1}{\epsilon}+1}\right)I^{\frac{1}{\epsilon}+1}\label{eq:total_cost_structure},
\end{equation*}

\noindent where $c_1$ scales up or down the cost of producing new housing and $\epsilon$ represents the price elasticity of housing investment. The problem of the construction company is to choose investment to maximize profits:

\begin{equation*}
\label{eq:construction_company_profits}
\underset{I} {\text{$\mbox{max} $}}{\left\{p I-TC(I)\right\}},
\end{equation*}

\noindent and the solution to this problem gives the following relationship between housing investment and the price of housing:

\begin{equation}
\label{eq:construction_company_first_order_condition}
I=\left(\frac{p}{c_{1}}\right)^{\epsilon}.
\end{equation}

\subsection{Aggregation}\label{subsection: aggregation}

We require that the government's budget balance in all our experiments and in the benchmark. The government's budget constraint,
\begin{align}\label{eq:gbc}
  & G + \sum_{j=R}^J \int_\mathcal{S} y_{R}(s) \mathrm{d}s + \sum_{j=1}^J \int_\mathcal{S} T \mathrm{d}s = \sum_{j=1}^J \int_\mathcal{S}y^{\tau}_{j}(s) \mathrm{d}s+\tau_p p H,
\end{align}

\noindent states that the government's expenditures on public goods, social security benefit payments, and lump sum transfers must equal to tax revenues from income and property. When we calibrate the model to the current tax code in Toronto we set the lump sum transfer $T$ to 0 and choose the consumption level of public goods, $G$, to ensure that this constraint holds. In our experiments, we keep $G$ constant, and adjust $T$ to clear the budget constraint. Note that $S=A\times O \times I \times H \times M \times E$ denotes a typical subset of the state space $\mathcal{S}$, $A$, $O$, $I$, $H$, $M$, and $E$, respectively.

The overall demand for housing must equal to the overall supply of housing:

\begin{equation}
\label{eq:housing_clearing_condition}
  & \sum_{j=1}^J \int_\mathcal{S} h_{j}(s) \mathrm{d}s = H',
\end{equation}

\noindent and the overall demand for rental housing must equal to the overall supply of rental housing:

\begin{equation}
\label{eq:rental_housing_clearing_condition}
  & \sum_{j=1}^J \int_\mathcal{S} (1-o'_{j}(s))h_{j}(s) \mathrm{d}s = X'.
\end{equation}

The distributions of households evolve according to the laws of motion,
\begin{align}
\Psi_{j+1}(S) &=\int_{\mathcal{S}} F(E,e)M_{j}(s)\mathrm{d}\Psi_{j}(s),\ j=1,\ldots,J-1\label{eq:lom1}
\end{align}
\begin{align}
\Psi_{1}(a=0,o=0,\iota=0,h=\underline{h},m=0,e) &=\int_{\mathcal{S}}\tilde{F}(E,e) \mathrm{d}\Psi_{J}(s),\label{eq:lom2}
\end{align}

\noindent where,

\begin{align}
    & M_{j}(s) = \indicator{[a'_{j}(s)\in A]) \land [o'_{j}(s)\in O] \land [\iota'_{j}(s)\in I] \land [h'_{j}(s) \in H] \land [m'_{j}(s) \in M]}.\notag
\end{align}

\subsection{Equilibrium}
Given an exogenous rate of decrease in real rents for every year a renter has lived in the same unit, $\kappa$, a stationary equilibrium in the model is a sequence of aggregate prices and quantities, $\left\{H',I,X',p,\overline{p}^{r}\right\}$, a sequence of value and policy functions, $\biggl\{V_{j}(\cdot),c_{j}(\cdot),a'_{j}(\cdot),o'_{j}(\cdot),\iota'_{j}(\cdot),h'_{j}(\cdot),m'_{j}(\cdot)\biggl\}_{j=1}^J$, and a sequence of measures, $\left\{\Psi_{j}(\cdot)\right\}_{j=1}^J$, that:
\begin{enumerate}
%\item all nominal prices rise at the rate of $i$, but real prices remain fixed;
\item solve the household's dynamic problem (\ref{eq:vf}) subject to constraints (\ref{eq:bc})-(\ref{eq:mortgage_debt_renter});
\item solve the rental company's problem by satisfying the rental price equation, (\ref{eq:rental_price});
\item satisfy the construction company's first-order condition, (\ref{eq:construction_company_first_order_condition});
\item satisfy the government's budget constraint, (\ref{eq:gbc});
\item clear the housing market,
(\ref{eq:housing_clearing_condition});
\item clear the rental housing market,
(\ref{eq:rental_housing_clearing_condition});
%\item satisfy the average declared house price as described in
%(\ref{eq:average_declared_price});
\item and satisfy the laws of motion for the distribution of households, (\ref{eq:lom1})--(\ref{eq:lom2}).
\end{enumerate}

\section{Calibration}

\subsection{Assigned parameter values}
Calibration follows two steps. First, we exogenously assign values to certain parameters that have been estimated in the literature. Second, we internally calibrate the values of the remaining parameters by picking their values so that the model generates certain statistics as observed in the data. Table \ref{tab:external_calibration} lists the parameter values, which we break into several groups.

We use 2016 census data of Toronto CMA to compute several key parameters. Homeownership rate is 67\%. Total rent to total income and renters with rent/income higher than 30\% (calculate in VDL). In 2016, the MLS average price is \$767,336 and average rent of primary apartments is \$1,240 from Rental Market Survey. The average rent of secondary apartments is \$1,901 and the average rent of other dwelling types other than apartments is \$1,583. Taking the average rent of primary and secondary markets, the ratio of rent over house price is 2.46\%. The HPI for apartments in 2016 is \$359,618 and average rent of apartments is \$1570.5. The ratio of rent over apartment price is 5.24\%. 
From CMHC (2018), the elasticity between housing starts and house prices in Toronto from 1992 to 2016 is 0.35. We also collected property tax rates from different municipalities in Toronto CMA and computed the average property tax rate, 0.67\%. The average real growth rate of house prices in Toronto from 2006 to 2016 is 4.46\% using MLS average house prices (that's more workable). The average growth rate of house prices in Toronto from 2006 to 2021 is 4.04\%. \textcolor{red}{I'll take the 4.46\% and see what happens!}. The ratio of total rent to total income before tax for renters in Toronto CMA in 2016 is 24.29\%. The share of renters with rent/income over 30\% is 45.16\% in Toronto in 2016. The average household employment income in Toronto in 2016 is \$74,722. 
Next steps:
\begin{enumerate}
    \item Progressivity of labor labor income taxation; run regression of \begin{equation}
        log(\mbox{after tax income})=log(a)+\psi log(\mbox{gross income})
    \end{equation} 
    After regression, constant = 0.6031, $\psi=0.9267$ with R-squared 0.98. 
    \item Average income tax rate for Toronto (we have for Canada from OECD but need for Ontario).
    I divided household income after tax by household total income for Toronto in 2016. I got 88.24\%, implying that the average tax rate is 11.76\%.  In the data, we have a group of households that do not pay tax or little. That made the average income tax rate quite low. Not sure whether our approach is correct. 1-household income after tax over household income. 
    \item abor productivity parameters
    \item Homeownership rate for each age 25-86 from the Census...
    In 2016, homeownership rates for 25-34 years 42\%, for 35-44 years 64.3\%, for 45-54 years 72.2\%, for 55-64 75\%, and for 65 years and over 75.1\%. SR: it is hard to compute for 65-74 and 75-86? Would be good to have some consistency.
    \item SR: I will compute the LTV for each age from 25-86
    \item Fraction of Toronto's rental market not under regulation?
    \item SR: I will compute wealth inequality from SFS.
    \item SR: What fraction of Toronto's population rented in 1991 relative to 2016? And how many households lived in Toronto in 1991 and in 2016? 
    \item SR: Lin, you'll see in the text there are some places where I put ?? below. Can you add some description of how and where you got the statistic from, for all these ??. 
    \item SR: need to think of retirement benefits.
\end{enumerate}

\subsection{Demographics and preferences}
Households start their lives at age 26, retire at age 66, and live to age 85, implying that lifespan $J=60$ and retirement age $R=41$. We set households' relative risk aversion to $\sigma=2$. The discount factor, $\beta$, the share of housing services in utility, $\gamma$, and utility advantage to homeownership, $\eta$, are determined in the second stage of our calibration procedure.

\subsection{Labor market productivity}
Labor market productivity of surviving households and inter-generational labor market productivity follow AR(1) processes,
\begin{equation*}
  \log e' = \rho_{e} \log e + \epsilon_{e},\quad \epsilon_{e}\sim N(0,\sigma_{e}),
\end{equation*}
\begin{equation*}
\log e^{\text{child}} = \hat{\rho}_{e} \log e^{\text{parent}} + \hat{\epsilon}_{e},\quad \hat{\epsilon}_{e}\sim N(0,\hat{\sigma}_{e}).
\end{equation*}
We set $\sigma_{e}=??$, $\rho_{e}=??$, $\hat{\sigma}_{e}=??$, and $\hat{\rho}_{e}=??$ to match the coefficients from intra- and inter-generational AR(1) regressions estimated using simulated data from the process for labor income in \citet{gkkoc}. We set the deterministic life-cycle component of labor market productivity to $\zeta_j =1+\min\{??,\frac{j}{R-11}??\}$ so that average households' labor income rises by ?? percent by age 55 \citep{guvenen-et-al}. We approximate retirement benefits, $B(e)$, using Canada's \citet{pspp}.

\subsection{Taxes and interest rate}
We set the consumption tax rate to 13\%, which is the Harmonized Sales Tax (HST) in Ontario. We set the rate at which capital income is taxable, $
\nu$, to 50\% as in Canada's tax code. We set the progressivity of the income tax system $\psi$ to ?? based on ...., and choose $\tau_{\ell}$, which determines the average level of taxation, in the second stage of our calibration. The interest rate, $r$, is set to 3\% as in \citet{kaplanetal-2019}. 

\subsection{Housing}

We pick the rate at which real rents decrease for renters remaining in the same unit, $\kappa$, to equal to 2.37\%, which we estimate using Census data from 2001-2016 and Rental Market Survey (RMS) data for 2016-2021. The exogenous real growth rate of house prices, $g_{p}$, that appears in the rental company's problem, is set to 4.46\% based on estimates from Multiple Listings Service (MLS) data from 2006-2016. We set the maximum ratio of mortgage debt to house value, $\lambda_{h}(ph')$, such that homeowners with a house that has a model-equivalent value above \$1M can borrow only up to 80\% of its value. Owners purchasing a home with a model-equivalent value between \$500,000-\$1M can borrow up to 90\% of the house value, and owners with a house value below \$500,000 can get a loan of up to 95\% of the house value. This is consistent with the rules and regulations in Canada's banking system. The interest rate premium on mortgage debt, $\Delta r_{m}$, is set to 1\% as in \citet{kaplanetal-2019}. The property tax rate, $\tau_p$, is set to 0.67\% based on ??, and the housing depreciation rate is chosen to equal 2\% as estimated by \citet{hardingetal-2007}. Buying costs, $\tau_{b}$, are set to 2.5\% of the purchase price and selling costs, $\tau_{s}$, are set to 7\% based on estimates by \citet{gm-2003}. Lastly, we choose $\epsilon$ to equal 0.35 to align with the estimate by ?? for the elasticity of the housing supply in Toronto.

\subsection{Internally calibrated parameters}

We now turn to internally calibrating our model. After assigning the parameter values listed above, there are seven parameters that still must be calibrated: the rate of time preference, $\beta$; the utility weight on the consumption of housing services, $\gamma$; the utility advantage to homeownership, $\eta$; the average level of income taxation, $\tau_{\ell}$; the convexity of managing rental housing, $\theta$; the average cost of managing rental housing in the controlled rental market, $c_{m,0}$; rent in the free rental market, $p^{r}_{1}$, and the smallest housing unit, $\underline{h}$.

We jointly calibrate all these parameters to match the following moments of Toronto data: the percentage of mortgage debt relative to gross housing wealth, which is 24\% on average according to the 2016 wave of the Survey of Financial  Security (SFS). The percentage of income renters spend on rent, which is 24.3\% and comes from 2016 Census. The homeownership rate, which is equal to 67\% according to the 2016 Census. The average income tax rate, which is ?? and comes from estimates by the ??. The tax incidence of renters, which is ?? and comes from estimates by ??. The rent to housing price ratio, which comes from ?? and is equal to 5.24\%. The share of the free rental market of the total rental market, which is ?? and comes from ??. The percentage of renters spending over 30\% of their gross income on rent, which comes from the 2016 Census and is equal to 45\%.

Roughly speaking, $\beta$ controls the mortgage debt to gross housing value ratio, $\gamma$ determines mostly the fraction of income renters spend on rent, $\eta$ affects the homeownership rate, $\tau_{\ell}$ determines the average income tax rate, $\theta$ determines the tax incidence of renters, $c_{m,0}$ determines the rent to housing price ratio, $p^{r}_{1}$ determines the share of the free rental market, and $\underline{h}$ controls the percentage of renters spending over 30\% of their gross income on rent.

\section{Non-targeted statistics}

\section{Deregulating rents}

\section{Conclusion}

\bibliographystyle{aer}
\bibliography{bib}
\linespread{1.5}

\input{calibration}
\input{non_targeted_statistics.tex}
\input{aggregate_results.tex}
\begin{table}[p]
\footnotesize
\renewcommand{\arraystretch}{1.2}
\begin{center}

\begin{tabular}{lccccccc}
\toprule
Quintile & Lowest & Second & Third & Fourth & Highest\\ 
\midrule
Unit control & -0.06 & -0.05 & -0.07 & -0.08 & 
-0.05 \\
Full decontrol & -0.28 & -0.32 & -0.36 & -0.36 & 
-0.23 \\
Full tenancy control & -0.03 & -0.02 & 0.00 & -0.11 & -0.06 \\
Welfare-maximizing degree of tenancy control & 0.01 & 0.01 & 0.01 & 0.01 & 0.01\\
\bottomrule
\end{tabular}
\end{center}
\normalsize
\end{table}


\end{document}